\documentclass[11pt]{article}
\usepackage[T1]{fontenc}
\usepackage[utf8]{inputenc}
\usepackage[a4paper,margin=2.4cm]{geometry}
\usepackage{amsmath,amssymb}
\usepackage{array,tabularx}
\pagestyle{plain}
\setlength{\parskip}{4pt}
\setlength{\emergencystretch}{2em}

\newcommand{\su}{\mathfrak{su}}

\begin{document}

\begin{center}
{\Large\bfseries Retraction of Theorems 1.1 and 3.11, and Erratum (v2)}\\[5pt]
{\large Gradient Flow Observables and Ultraviolet Closure without\\
Blocking-Map Hypotheses for Lattice Yang--Mills Theory in Four Dimensions}\\[6pt]
Lluis Eriksson \quad\textbullet\quad ai.viXra:2602.0085 \quad\textbullet\quad July 2026
\end{center}

\noindent\rule{\textwidth}{0.8pt}

\vspace{2pt}
\noindent
\textbf{The principal results of this paper are withdrawn.} Lemma 3.6 is false as
stated, and with it Lemma 3.8, Proposition 3.9, Theorem 3.11 and Theorem 1.1. The
defect is not a gap to be filled: the operator identity on which the whole
ultraviolet estimate rests is contradicted at the simplest configuration on the
lattice. This note states the refutation, delimits exactly what survives, and
records the repair route. The v1 manuscript follows unmodified, so that the claim
and its refutation can be read against each other.

\section{Status}

\begin{center}
\begin{tabularx}{\textwidth}{|>{\raggedright\arraybackslash}p{0.43\textwidth}|X|}
\hline
\textbf{Statement} & \textbf{Status as of v2}\\
\hline
Lemma 2.2 (gauge covariance of the flow) & Stands\\
\hline
Lemma 3.2 (Gaussian bound for the scalar kernel $p_\tau$)
 & \textbf{Recoverable} as an independently\\
 & specified scalar heat-kernel statement;\\
 & the original instantiation is \textbf{withdrawn}\\
\hline
Proposition 1.3 (flow--reflection commutation) & Stands\\
\hline
Lemma 3.6, Eq.~(16) (structure of $\mathcal L_\tau$) & \textbf{FALSE} (\S2)\\
\hline
Lemma 3.8 (domination by scalar heat semigroup) & \textbf{Withdrawn}\\
\hline
Proposition 3.9 ($\ell^2$ column bound) & \textbf{Withdrawn}\\
\hline
Theorem 3.11 (squared-oscillation summability) & \textbf{Withdrawn}\\
\hline
Theorem 1.1 (UV closure) & \textbf{Withdrawn}\\
\hline
Reflection positivity, OS, thermodynamic limit, mass gap & Open (as in v1)\\
\hline
\end{tabularx}
\end{center}

\section{The refutation}

Equation~(16) of Lemma 3.6 asserts that the linearization of the Wilson flow, in
transported coordinates, has the form
\[
\mathcal L_\tau \;=\; \Delta^{\mathrm{cov}}_{V_\tau} + A_{V_\tau},
\qquad
(\Delta^{\mathrm{cov}}_{V_\tau}X)(e)=\sum_{e'\sim e} a_{e,e'}\bigl[T_{e,e'}(\tau)X(e')-X(e)\bigr],
\]
with weights $a_{e,e'}\ge 0$ \emph{depending only on the lattice geometry}, $T_{e,e'}(\tau)$
fibre isometries, and $A_{V_\tau}$ acting pointwise by adjoint commutators. Lemma 3.8
and Proposition 3.9 use nothing about the flow except this structure.

\paragraph{The test configuration.}
Take $U\equiv \mathbf 1$ on the torus. Every plaquette is trivial, so $U$ is a global
minimum of the Wilson action $S_W$ and a fixed point of the flow: $V_\tau\equiv\mathbf 1$
for all $\tau$. At this configuration all parallel transports are trivial,
$T_{e,e'}(\tau)=\mathrm{id}$, and the plaquette holonomies entering $A_{V_\tau}$ are
trivial, so $A_{V_\tau}=0$. Equation~(16) therefore asserts that $\mathcal L_\tau$ is
\emph{exactly} a graph Laplacian with non-negative weights,
$(\mathcal L X)(e)=\sum_{e'\sim e}a_{e,e'}[X(e')-X(e)]$.

\paragraph{A dimension count that (16) cannot survive.}
Let $\mathcal G_+$ be the subgraph of the link graph formed by those pairs $e\sim e'$
with \emph{strictly} positive weight $a_{e,e'}>0$. The two cases below are exhaustive,
and the chain of \S3.3 breaks in each.

\emph{Case 2: $\mathcal G_+$ is disconnected.} Then the scalar comparison semigroup
$P_\tau=e^{\tau\Delta}$, which by construction carries the \emph{same} weights, is
reducible: $p_\tau(\cdot,e_0)$ remains supported in the connected component
$C(e_0)\subsetneq\mathcal G_k$, and $p_{2\tau}(e_0,e_0)\to 1/|C(e_0)|>1/|\mathcal G_k|$
as $\tau\to\infty$. The scalar-kernel bound invoked in the proof of Proposition 3.9,
$\sum_e p_\tau(e,e_0)^2=p_{2\tau}(e_0,e_0)\le C_{\mathrm{jac}}(\tau+1)^{-2}
+1/|\mathcal G_k|$, is then false for $\tau$ large. The argument fails before the gauge
modes are reached: the single global stationary term $1/|\mathcal G_k|$ presupposes
irreducibility.

\emph{Case 1: $\mathcal G_+$ is connected.} Then the kernel of the claimed weighted
Laplacian is exactly the space of constant $\su(N)$-valued functions, of dimension
$\dim\su(N)=N^2-1$ --- and this is contradicted by the gauge count that follows.

The true linearization at a fixed point of a gradient flow is minus the Hessian:
$\mathcal L=-\operatorname{Hess}S_W(\mathbf 1)$. Now $S_W$ is gauge invariant, and the
gauge orbit of $\mathbf 1$ consists entirely of pure-gauge configurations
$U^g(e)=g(x)g(y)^{-1}$, all of whose plaquettes are trivial; hence $S_W$ is
\emph{constant, at its minimum,} along the whole orbit, which is therefore a manifold
of minima. (Constancy and minimality are all that is used; whether the constant is
zero depends on the normalization chosen for the action.) The Hessian annihilates the
tangent space to a manifold of minima, so
\[
\mathcal L\,X^\lambda=0
\qquad\text{for every }
X^\lambda(e):=\lambda(x)-\lambda(y),\quad e=(x,y),\ \ \lambda:\Lambda^0_k\to\su(N).
\]
The gauge-tangent subspace $\{X^\lambda\}$ has dimension $(|\Lambda^0_k|-1)(N^2-1)$
(the constant $\lambda$ act trivially). Only a \emph{lower} bound is needed, and only
a lower bound is claimed --- $\mathcal L$ may well have further null modes, and none
need be excluded:
\[
\dim\ker\mathcal L\;\ge\;(|\Lambda^0_k|-1)(N^2-1)\;>\;N^2-1
\qquad\text{whenever }|\Lambda^0_k|>2,
\]
in particular on every lattice of the refinement sequence considered here (the
inequality is an equality, not a strict one, at exactly two vertices).
whereas the claimed connected weighted Laplacian has kernel of dimension
\emph{exactly} $N^2-1$. \textbf{The two conclusions about $\dim\ker\mathcal L$ are
incompatible.} Equation~(16) is false, and it is false already at $U\equiv\mathbf 1$,
with no asymptotics and no field-dependence involved.

So in Case 1 the asserted operator identity is refuted outright, and in Case 2 the
heat-kernel estimate built on it is refuted outright. No reading of ``weighted
connected graph'' rescues \S3.3.

\paragraph{The same defect, measured against Proposition 3.9.}
Take $\lambda$ supported at a single vertex $x_0$ with $|\lambda(x_0)|=1$. Then $X^\lambda$
is supported on the $2d$ links incident to $x_0$, with $|X^\lambda(e)|=1$ on each, and by
the above it is \emph{exactly stationary}: $X^\lambda_\tau=X^\lambda_0$ and
$\sum_e|X^\lambda_\tau(e)|^2=2d$ for all $\tau$. But $X^\lambda$ is a signed sum of $2d$
single-link initial data, so linearity of~(15) and the triangle inequality in
$\ell^2(\mathcal G_k)$ give, from Proposition 3.9,
\[
\sqrt{2d}\;=\;\bigl\|X^\lambda_\tau\bigr\|_{\ell^2}
\;\le\; 2d\,\sqrt{\frac{C_{\mathrm{jac}}}{(\tau+1)^2}+\frac{1}{|\mathcal G_k|}}\,.
\]
Choosing $\tau$ large enough that $C_{\mathrm{jac}}(\tau+1)^{-2}\le|\mathcal G_k|^{-1}$
yields $|\mathcal G_k|\le 4d$. For $d=4$ every torus with more than $16$ links
contradicts the proposition. The $1/|\mathcal G_k|$ term that Remark 3.10 dismisses as
subdominant is precisely where the gauge modes were lost.

\paragraph{Where the proof of Lemma 3.6 goes wrong.}
Differentiating the staple products in $\partial_{x,\mu}S_W$ produces, besides the
nearest-neighbour transport--difference terms and the on-site adjoint commutators, a
term of divergence type --- in the continuum linearization
$\partial_\tau b_\mu=D^2b_\mu+2[G_{\mu\nu},b_\nu]-D_\mu(D_\nu b_\nu)$, the summand
$-D_\mu(D_\nu b_\nu)$. It is second order, it is neither a non-negative-weight
transport--difference nor a pointwise adjoint commutator, and it is exactly the term
that annihilates longitudinal (gauge) modes. The proof of Lemma 3.6 sorts the
differentiated terms into two buckets and has no bucket for it. This is the same
object that L\"uscher and Weisz damp by hand with the parameter $\alpha_0$: in the
linearized flow the transverse part carries $e^{-\tau p^2}$ and the longitudinal part
$e^{-\alpha_0\tau p^2}$, so at $\alpha_0=0$ --- the flow used in this paper --- the
longitudinal part does not decay at all.

\section{What this costs, and what it does not}

Theorem 3.11 has no other proof in the paper, and Theorem 1.1 obtains its influence
bound (\S4.4, Proposition 4.4) from Theorem 3.11. Both are withdrawn. In particular
the paper's central assertion --- that the oscillation summability which was an
\emph{assumption} on the blocking map in the companion paper is here a \emph{theorem}
--- is retracted; that assumption is not discharged by this work.

Two results of the author stand exactly as printed, and a third scalar estimate is
recoverable but not standing; none of them uses (16):
\begin{itemize}
\item \textbf{Lemma 2.2}, gauge covariance of the flow. It not only survives, it is the
instrument of the refutation above.
\item \textbf{Lemma 3.2}, the Gaussian upper bound for the scalar heat kernel $p_\tau$
on the link graph --- but only in a restricted sense, and the restriction matters. As
a statement about a graph Laplacian that is \emph{specified independently}, connected
and uniformly elliptic, it is standard and survives; it has no gauge content. What
does \emph{not} survive is its role in this paper: the weights defining $p_\tau$ were
taken from the decomposition (16), which is false, so $p_\tau$ is no longer the
dominating semigroup of anything, and in the disconnected branch of \S2 the bound with
the single stationary term $1/|\mathcal G_k|$ is itself false. Lemma 3.2 has no
remaining implication for the Wilson-flow variational equation.
\item \textbf{Proposition 1.3}, $W_\tau\circ\Theta=\Theta\circ W_\tau$. Proved from
invariance of the action under reflection and uniqueness for the flow ODE. It is
elementary, and the paper already states that it does not by itself give reflection
positivity for the smoothed observables.
\end{itemize}

\noindent
The proofs of the two surviving results of the author were re-read for this note,
rather than cleared by inspecting what they depend on: Lemma 2.2 follows from gauge
invariance of $S_W$, covariance of the link derivative under gauge transformations,
and uniqueness for the flow ODE; Proposition 1.3 from $\Theta$-symmetry of the action,
$S_W(\Theta U)=S_W(U)$, and the same uniqueness. Both are sound. The distinction is
not pedantry: in the companion retraction (ai.viXra:2602.0084 v2) a statement cleared
the other way --- by checking its dependencies without reading its proof --- turned
out to be false as printed.

\paragraph{A dangling forward pointer.} Remark 5.1 of the manuscript refers the reader
to a companion paper in which ``Strategy (a)'' --- approximation by
half-space-supported observables via conditional expectation, with Gaussian
localization controlling the $L^2$ error uniformly in $k$ --- ``is carried out in
detail''. That companion is \textbf{ai.viXra:2602.0084}, whose corresponding results
are withdrawn in its own v2, in part \emph{because} they depend on Theorem 1.1 of the
present paper. The two manuscripts cite each other across this gap, and neither
supplies what the other assumes. The forward pointer of Remark 5.1 is open, not
delivered.

\emph{And it is more open than this paper's defect alone would imply.} What the
present retraction propagates to that companion is one defect: the non-existence of
the limit state. Its own v2 records four further defects that owe nothing to the
gradient flow --- in its variance--oscillation estimate, in the periodic support
geometry of the torus, in a chain rule applied to a nonlinear map, and in its
finite-volume heat-kernel input. Repairing the present paper would therefore not
repair that one; the damped-flow and orbit-space routes of \S4 below address only the
defect propagated from here.

\paragraph{Effect on companion manuscripts and on the formalized lane.}
No compiled theorem in the Lean core of the repository
\texttt{THE-ERIKSSON-PROGRAMME} discharges a hypothesis of the present paper, and none
is affected by this retraction; the formalized lane and this paper do not meet.
At the level of manuscripts the situation is not empty, and is stated here so that it
is not discovered later:
\begin{itemize}
\item \textbf{ai.viXra:2602.0084} (\emph{Almost Reflection Positivity for
Gradient-Flow Observables via Gaussian Localization}) imports the withdrawn Lemma 3.6
by citation, as its own Lemma 3.1, in exactly the form refuted here (a bound
$\|X_\tau(\ell,\ell_0)\|_{\mathrm{op}}\le p_\tau(\ell,\ell_0)$ on the Jacobian
columns). Its Gaussian-localization conclusion is therefore \emph{not presently
established}, and must be stated conditionally on a corrected damped or transverse
estimate. \textbf{This retraction does not by itself prove that that conclusion is
false}: its consumers there require gauge-invariant observables, for which repair
route (b) below remains open.
\item \textbf{ai.viXra:2602.0077} is unaffected. Its summability input is an explicit
\emph{hypothesis} concerning the blocking map $Q_{\ell,k}$, not a theorem about the
Wilson flow, and that paper's own v2 had already downgraded its discharge statuses.
\end{itemize}

\section{Repair route (a route, not a result)}

The obstruction is specific: the unfixed Wilson flow has undamped gauge directions, so
\emph{no} bound of the form ``every Jacobian column decays'' can hold. Two observations
make a repair plausible, and neither is carried out here.

\begin{enumerate}
\item[(a)] For a \emph{gauge-invariant} $F$, the composed observables coincide as
functions on configuration space, $F\circ W^{(\alpha_0)}_\tau=F\circ W^{(0)}_\tau$,
because the two flows differ by a $\tau$-dependent gauge transformation. Hence
$\operatorname{osc}_e$ of the two is literally equal, and it suffices to prove the
oscillation bound for the damped flow $\alpha_0>0$, whose linearization is uniformly
parabolic in \emph{all} directions. The structure (16) is then not obviously wrong ---
but it must be re-derived, with the damping term included and its sign controlled.
\item[(b)] Alternatively, work on the orbit space and estimate only transverse to the
gauge directions. This requires a Jacobian estimate on a quotient, not on
$\su(N)^{\mathcal G_k}$, and the counting above shows the two are genuinely different
problems.
\end{enumerate}

\noindent
Either route needs a new proof of the oscillation bound. Until one exists, the
ultraviolet closure of this paper stands at the same place as the companion paper's:
conditional on an unproved summability input.

\section{Erratum: the archive metadata of v1}

Separately from the mathematics, the v1 archive record carried the title and abstract
of an \emph{earlier draft}, claiming ``\emph{Wilson-loop expectations \ldots\ with
compact gauge group $G$ admit a universal continuum limit}''. Even before the
retraction above, Theorem 1.1 was narrower on five axes: fixed finite torus (the
thermodynamic limit is Open, \S6.4); gradient-flow observable families at fixed
physical flow time $t>0$, not Wilson loops (the limit $t\downarrow 0$ is never taken);
$\mathrm{SU}(N)$ rather than general compact $G$; weak coupling $g_0\in(0,g_*]$ only;
and conditional on Assumption A together with the hypotheses (H1)--(H2) imported in
\S4.3. The manuscript itself never made the stronger claim. The record's title and
abstract have been corrected.

Two further overstatements internal to v1 are withdrawn, and would have been withdrawn
even without the refutation: the abstract's description of Assumption A as ``verified
automatically for all standard physical observables'', and the corresponding row of the
\S6.7 table. Remark 4.3 is titled ``\emph{Expected} verification of Assumption A'' and
states that the estimate ``is \emph{expected} to hold''; it sketches a two-region
mechanism without explicit constants. The label ``unconditional'' is withdrawn from the
title of \S4 and from \S6.7.

\vspace{4pt}
\noindent\rule{\textwidth}{0.8pt}

\noindent\footnotesize
\textbf{Provenance and credit.} The metadata defect was measured on 2026-07-29 by
comparing the archive metadata of every record in the author's viXra listing against
the cover page of the served PDF (93 records; 5 genuine mismatches). \textbf{The
refutation in \S2 originates in an external adversarial review of this paper}, which
identified the undamped gauge directions of the unfixed Wilson flow as incompatible
with Lemma 3.6 and cited L\"uscher--Weisz (arXiv:1101.0963) for the role of $\alpha_0$;
the reviewer's argument was verified line by line against the manuscript before this
note was written, and the kernel-dimension form of the contradiction given in \S2 was
obtained in that verification. A second round of the same review supplied the case
split on $\mathcal G_+$ in \S2, without which the dimension count would have rested on
an unstated connectivity assumption, and required the lower-bound form of
$\dim\ker\mathcal L$. The error was the author's; the catch was not.

\end{document}
